\section{算力与组织约束：开放科学的现实瓶颈}

\subsection{Compute 仍是第一性约束}
讲者指出，即便开放定义清晰、社区意愿强烈，开放科学仍会被算力与系统工程约束。  
这意味着未来竞争不只在模型能力，也在谁能构建可复用的训练与评测基础设施。

\begin{figure}[H]
\centering
\includegraphics[width=0.93\textwidth]{frame-07-adoption.jpg}
\caption{生态采用与社区贡献：开放框架的网络效应}
\end{figure}
\noindent\footnotesize 画面证据：字幕约在 00:26:01 与 00:31:01，讲者讨论生态规模、开放模型实践与研究迁移价值。\normalsize

\begin{knowledgebox}{基础设施视角下的三类稀缺资源}
\begin{itemize}
    \item 训练算力：决定可探索的模型/数据空间；
    \item 工程人才：决定系统可维护性与实验吞吐；
    \item 评测资产：决定进展是否可验证与可比较。
\end{itemize}
\end{knowledgebox}

\subsection{去中心化协作的可行路径}
\begin{importantbox}{开放生态的真正杠杆：共享标准与共享工具链}
单点算力难以复制 frontier 训练，但社区可以通过统一评测、可移植训练脚本、可复用数据治理流程形成“协作乘数”。
\end{importantbox}

\begin{warningbox}{没有标准化，协作会退化成碎片化重复劳动}
若各团队使用互不兼容的日志格式、评测协议和任务定义，结果是“每个人都在做实验，但没人能有效比较实验”。
\end{warningbox}

\subsection{本章小结}
开放科学不等于低门槛。它需要新的共同基础设施来降低协作摩擦，否则“开放”只会停留在口号层面。


